\section{Limitations}
\label{sec:limits}

\textbf{Compute-capped opponents.} Each opponent run had a fixed wall-clock budget on one T4 (30--55 minutes for
STAEformer and Graph WaveNet, 6--12 for STID; Table~\ref{tab:budget}). STID reached its authors' 100-epoch schedule
on four networks; Graph WaveNet on one; STAEformer on none, and it ran out of memory on PEMS07 at its published batch
size. The comparison is therefore compute-matched rather than converged: STAEformer's numbers in particular
understate what it reaches with its full schedule, and its published results (e.g.\ 13.46 MAE on PEMS08 with all test
windows) remain ahead of ours where reported. Each opponent was run with one seed.

\textbf{Development exposure.} PEMS08 and PEMS04 shaped PWM-Dyn's design in earlier versions, including one look at
PEMS08 test scores of development checkpoints; the reported models use fresh seeds and a configuration chosen on
validation, but the design itself is not blind on those two networks. PEMS03, PEMS07, METR-LA and PEMS-BAY were never
used during development.

\textbf{Scope.} Traffic tasks use the 12-step literature protocol and the new tasks score every sixth window.
Rainfall covers two Indian regions and its PWM-Dyn entries come from the previous model version. TimesFM 3.0 is used
zero-shot under a non-commercial research licence; we did not fine-tune it, and fine-tuned foundation models may
close the gap we report. A PWM-Dyn win over a foundation model is a specialist-versus-zero-shot result, not evidence
of better generalisation; the transfer track shows the opposite.

\section{Conclusion}

Trained on a network's own history and scored on identical windows, a $3\times10^5$-parameter graph model with
explicit missing-data inputs and conformal intervals is the most accurate of the trained models we ran on five of six traffic networks, ties the best fully trained
specialist on two of them, and beats a 2026 zero-shot foundation model on most tasks in both point and interval
quality, while degrading gracefully when sensors fail. It is not state of the art in absolute terms, loses to
STAEformer on PEMS08 and to TimesFM~3.0 on PEMS03, and its advantage over foundation models depends on having the
network's own history.
 The benchmark, its scorer, every
opponent harness and all saved forecasts are released so that these comparisons can be re-run and extended.
